% contexto del problema,
% motivación,
% hipótesis o idea general del trabajo,
% objetivos generales y específicos,
% metodología general del trabajo,
% estructura de la memoria.

\chapter{Introducción}
\label{chap:introduccion}

La optimización de problemas complejos continuos constituye un área de alta relevancia dentro de la informática y, en particular, dentro del marco de la inteligencia artificial y la computación evolutiva. En un gran número de aplicaciones reales, el objetivo consiste en encontrar configuraciones de alta calidad dentro de espacios de búsqueda de gran dimensión y, normalmente, difíciles de abordar mediante métodos exactos o deterministas. La dificultad no reside únicamente en el tamaño del espacio de búsqueda, sino también en la complejidad del propio problema de optimización.

En este contexto, las metaheurísticas han adquirido un papel relevante debido a su capacidad para explorar espacios de búsqueda complejos. A diferencia de otros métodos de optimización, estas técnicas permiten trabajar con problemas multimodales, no lineales o de estructura desconocida, lo que explica su utilización cuando no existen métodos exactos adecuados o cuando su aplicación no es óptima computacionalmente. Entre ellas, los algoritmos evolutivos destacan por combinar mecanismos de exploración y explotación para aproximar soluciones de alta calidad en problemas de muy diversa naturaleza.

Sin embargo, esta versatilidad tiene un coste importante: las metaheurísticas suelen necesitar muchas evaluaciones de la función objetivo para obtener soluciones competitivas. En la mayoría de algoritmos evolutivos, cada solución candidata debe evaluarse para estimar su calidad y orientar el proceso de búsqueda. Cuando dicha evaluación es sencilla, este proceso no supone un problema computacional. No obstante, en problemas reales donde la función objetivo representa simulaciones complejas, procesos físicos o procedimientos numéricos de alta carga computacional, el tiempo total de ejecución queda dominado por el coste acumulado de esas evaluaciones. En estos casos, incluso metaheurísticas bien diseñadas resultan inviables si no se dispone de mecanismos que reduzcan dicho esfuerzo computacional.

Esta limitación ha impulsado el interés por investigar estrategias que permitan
mantener la capacidad exploratoria de las metaheurísticas reduciendo el número
de evaluaciones reales. Entre las alternativas más utilizadas se encuentran las
técnicas subrogadas, capaces de estimar el comportamiento de la función objetivo
a partir de los datos observados durante el propio proceso de optimización. De
este modo, la información obtenida por la metaheurística puede utilizarse no
solo para actualizar la población, sino también para construir un modelo
auxiliar capaz de orientar la búsqueda de forma más eficiente. En este tipo de
integración, la utilidad del subrogado no depende únicamente de aproximar con
precisión el valor numérico de la función objetivo, sino también de preservar
correctamente las relaciones de calidad entre soluciones candidatas, ya que
muchas decisiones evolutivas se basan en comparar individuos más que en conocer
su valor exacto.

\section{Motivación}
\label{sec:motivacion}

Debido a que el aprendizaje automático y la computación evolutiva se han desarrollado tradicionalmente como dos comunidades científicas independientes, existe una falta de retroalimentación mutua en el diseño de algoritmos híbridos. La oportunidad de conectar ambas disciplinas motiva y define el objetivo de esta investigación. Sin embargo, integrar un modelo de aproximación en un procedimiento evolutivo plantea dificultades, ya que la distribución de las soluciones evaluadas  cambian conforme la población evoluciona y dependen de la naturaleza de la propia metaheurística. Esta característica hace que no baste con entrenar un modelo de forma aislada, sino que sea necesario fundamentar la hibridación en un análisis de los datos generados por la metaheurística y cómo pueden aprovecharse.

Con el objetivo de conectar estas dos comunidades, en este Trabajo de Fin de Grado vamos a estudiar y mejorar el proceso de cómo se pueden utilizar los algoritmos de aprendizaje automático a partir del histórico de soluciones de los distintos algoritmos evolutivos para aplicarlas como técnicas de subrogado, y vamos a demostrar experimentalmente cómo se comportan.
Para ello, evaluaremos el rendimiento de las metaheurísticas sobre un banco de problemas continuos a través de un enfoque metodológico dividido en dos etapas. En primer lugar, analizaremos el histórico de datos recolectado por los algoritmos para determinar con precisión qué técnicas y bajo qué configuraciones de hiperparámetros muestran mayor utilidad para ordenar soluciones candidatas. En segundo lugar, utilizaremos las conclusiones extraídas para integrar los mejores modelos detectados dentro de los algoritmos evolutivos, validando de manera dinámica su capacidad real para guiar eficazmente el proceso de optimización.

\section{Objetivos}

El objetivo principal de este Trabajo de Fin de Grado consiste en estudiar, implementar y evaluar el uso de modelos de aproximación como técnicas de subrogado para asistir a algoritmos evolutivos en problemas de optimización continua.

En resumen, los objetivos de este trabajo son:

\begin{enumerate}
    \item Evaluar el comportamiento de las metaheurísticas seleccionadas sobre el banco de problemas considerado, incorporando estrategias de reinicio orientadas a mitigar la convergencia prematura y a favorecer la obtención de datos de búsqueda más informativos.

    \item Construir y evaluar estrategias de entrenamiento y validación que respeten la estructura temporal de las muestras, con el fin de comparar distintas familias de modelos subrogados, incluyendo algoritmos de aprendizaje automático y aproximadores clásicos. Esta comparación se centra principalmente en la capacidad de ordenación de los modelos, considerando también su coste computacional.

    \item Diseñar, implementar y evaluar la integración de los modelos más prometedores dentro de los algoritmos evolutivos, comparando su impacto final sobre la calidad de las soluciones, la dinámica de convergencia y el ahorro efectivo en el uso de evaluaciones reales, frente a las metaheurísticas originales.
\end{enumerate}

\section{Estructura de la memoria}

El presente trabajo se organiza en varios capítulos que recogen de forma
progresiva el contexto del problema, la planificación del proyecto, los
fundamentos teóricos necesarios, la revisión de la literatura, el diseño
experimental, el análisis de resultados y las conclusiones derivadas del
estudio realizado.

En primer lugar, el Capítulo~\ref{chap:introduccion} introduce el problema
abordado, expone la motivación del trabajo y formula los objetivos perseguidos.
A continuación, el Capítulo~\ref{chap:planificacion} describe la planificación
seguida durante el desarrollo del proyecto, así como la estimación del coste
asociado a su realización.

El Capítulo~\ref{chap:fundamentos_teoricos} desarrolla los fundamentos
teóricos necesarios para situar el trabajo, incluyendo los conceptos básicos de
optimización, metaheurísticas, aprendizaje automático y modelos subrogados. El
Capítulo~\ref{chap:estado_arte} recoge el estado del arte, revisando los principales enfoques existentes en la literatura sobre la integración de modelos subrogados en procesos de optimización evolutiva.

Seguidamente, el Capítulo~\ref{chap:algo_tecnicas_ref} presenta las
metaheurísticas, los modelos subrogados y las estrategias de referencia
utilizadas durante el estudio. El
Capítulo~\ref{chap:metodologia_experimental} describe el diseño experimental
aplicado sobre dichas componentes, incluyendo el benchmark, las métricas, las
estrategias de evaluación, el ajuste de configuraciones y los detalles
principales de implementación.

El Capítulo~\ref{chap:experimentacion} recoge los resultados experimentales
obtenidos y analiza el comportamiento de las metaheurísticas, la evaluación de
los modelos subrogados y su posterior integración en el proceso de
optimización.

Por último, el Capítulo~\ref{chap:conclusion} resume las principales
conclusiones alcanzadas y plantea las líneas de trabajo futuro. La memoria se
completa con apéndices que recogen tablas complementarias utilizadas para
apoyar el análisis experimental, así como los parámetros necesarios para
reproducir las ejecuciones principales.

